\documentclass[12pt]{article}
\usepackage[utf8]{inputenc}
\usepackage[T1]{fontenc}
\usepackage{hyperref}
\title{Wajiara Language Resources for K-12 Curriculum Development}
\author{Compiled by Claude Sonnet 5.5 (claude-sonnet-5-5) for Andr\'as Kornai}
\date{}
\begin{document}
\maketitle

\section{Introduction}
Wajiara (Yurut\'i, Juriti; ISO 639-3 \texttt{yui}, Glottocode yuru1263) is an Eastern Tucanoan language. Wikipedia gives about 1,200 speakers (2011) along the Paca and Ca\~no T\'i rivers in Vaup\'es, Colombia, and in the Alto Rio Negro territory in Brazil. Sources disagree on its vitality. Endangered Languages cites a 2005 census count of 266 speakers in Colombia, and Joshua Project gives a larger figure. My searches found no school materials and very little published documentation.

\section{Primary Educational Resources}
I found no primers, textbooks or school programmes for Wajiara.

\section{Linguistic and Reference Materials}
\begin{itemize}
\item Wikipedia, ``Wajiara language'': overview and phoneme inventory. \url{https://en.wikipedia.org/wiki/Wajiara_language}
\item Glottolog entry, with alternate names and references. \url{https://glottolog.org/resource/languoid/id/yuru1263}
\item SAPhon (Berkeley) phonological inventory for Yurut\'i; it notes nasal harmony and tone. \url{https://linguistics.berkeley.edu/saphon/en/inv/Yuruti.html}
\item Yurut\'i document hosted at dice.missouri.edu. I did not confirm its contents. \url{https://dice.missouri.edu/assets/docs/tukano/Yuruti.pdf}
\item Colombian Ministry of Culture document on the Yurut\'i people. I did not confirm its contents. \url{https://www.mincultura.gov.co:443/areas/poblaciones/noticias/Documents/Yurut\%C3\%AD.pdf}
\item Numerals page for Wajiara (MPI-EVA Channumerals). \url{https://lingweb.eva.mpg.de/channumerals/Wajiara.htm}
\item Omniglot page on Yurut\'i. \url{https://www.omniglot.com/writing/yuruti.htm}
\end{itemize}
The search results name Kinch and Kinch (2000), ``El yurut\'i'', as the main descriptive source. I did not find a URL for it.

\section{Religious and Community Resources}
I found no scripture or community organisation resources. Joshua Project has a population profile. \url{https://joshuaproject.net/index.php/people_groups/16076/CO}

\section{Digital and Audio Resources}
Endangered Languages has a catalogue entry. \url{https://endangeredlanguages.com/lang/2770}
I found no audio or video resources.

\section{Curriculum Development Context}
Documentation is thin, with no known dictionary, grammar or school programme. A curriculum would need new community-led documentation, working with Yurut\'i speakers and the relevant Colombian and Brazilian indigenous education authorities. Kinch and Kinch (2000) is the best starting reference.

\end{document}